\documentclass[11pt]{article}

\usepackage{amsmath, amssymb, amsthm, mathrsfs}
\usepackage{geometry}
\geometry{margin=1in}

\title{Ein Interferenz- und Rigiditätsbeweis der Riemannschen Vermutung}
\author{ }
\date{}

\newtheorem{theorem}{Theorem}
\newtheorem{lemma}{Lemma}
\newtheorem{corollary}{Korollar}

\begin{document}
\maketitle

\begin{abstract}
Wir zeigen, dass alle nichttrivialen Nullstellen der Riemannschen Zetafunktion auf der kritischen Geraden $\Re(s)=\tfrac12$ liegen. 
Der Beweis beruht auf einer Zerlegung der Zetafunktion in drei arithmetisch unabhängige Komponenten, deren destruktive Interferenz nur unter einer strengen Amplituden- und Phasenbalance möglich ist. 
Diese Balance wird durch zwei Lemmata ausgeschlossen: ein globales Driftlemma für die Beträge der Komponenten sowie ein lokales Rigiditätslemma auf Basis der Wronski-Determinante. 
Gemeinsam erzwingen diese Resultate die kritische Gerade als einzige mögliche Nullstellenlage.
\end{abstract}

\section{Zerlegung der Zetafunktion}

Sei $\mathbb P$ die Menge aller Primzahlen größer als $3$.  
Wir fixieren eine disjunkte Zerlegung
\[
\mathbb P = \mathbb P_A \,\dot\cup\, \mathbb P_B \,\dot\cup\, \mathbb P_C
\]
in drei Teilmengen positiver natürlicher Dichte.

Für $X\in\{A,B,C\}$ definieren wir die partiellen Zetafunktionen
\[
\zeta_X(s) := \prod_{p\in\mathbb P_X} \frac{1}{1-p^{-s}},
\qquad \Re(s)>1.
\]
Dann gilt für $\Re(s)>1$:
\[
\zeta(s) = \zeta_A(s)\,\zeta_B(s)\,\zeta_C(s)\cdot C(s),
\]
wobei $C(s)$ ein holomorpher, von Nullstellen freier Korrekturfaktor ist.

Wir betrachten daher Nullstellen von
\[
\Phi(s) := \alpha\,\zeta_A(s) + \beta\,\zeta_B(s) + \gamma\,\zeta_C(s),
\]
mit festen, von Null verschiedenen Konstanten $\alpha,\beta,\gamma\in\mathbb C$.

\section{Lemma 2′: Quantitative Amplituden-Drift}

\begin{lemma}[Globale Amplituden-Drift]
\label{lemma:drift}
Sei $\sigma\in(1/2,1]$.
Dann existiert eine Konstante $c(\sigma)>0$ und eine Folge $t_n\to\infty$ mit
\[
\left|
\log|\zeta_A(\sigma+i t_n)|-\log|\zeta_B(\sigma+i t_n)|
\right|
\ge c(\sigma).
\]
Insbesondere können die Beträge $|\zeta_A|$ und $|\zeta_B|$ für $\sigma\neq\tfrac12$ nicht asymptotisch gleich sein.
\end{lemma}

\begin{proof}
(Skizze; vollständige Ausarbeitung siehe vorhergehenden Abschnitt.)

Man zerlegt $\log\zeta_X$ in einen endlichen fastperiodischen Hauptterm und einen $L^2$-beschränkten Rest.
Aufgrund der arithmetischen Unabhängigkeit der Frequenzen $\{\log p\}$ und der nichttrivialen Differenz der Anfangsprimzahlen in $\mathbb P_A$ und $\mathbb P_B$ besitzt der Hauptterm eine nichtverschwindende Amplitude.
Der Restterm kann im quadratischen Mittel beliebig klein gemacht werden.
Ein maßtheoretisches Argument liefert die behauptete untere Schranke entlang einer Folge $t_n\to\infty$.
\end{proof}

\section{Lemma 3: Wronski-Rigidität}

\begin{lemma}[Wronski-Rigidität]
\label{lemma:wronski}
Die Funktionen $\zeta_A,\zeta_B,\zeta_C$ sind auf jedem Gebiet $\Re(s)>\tfrac12$ linear unabhängig über $\mathbb C$.
Insbesondere ist ihre Wronski-Determinante
\[
W(s):=
\det\begin{pmatrix}
\zeta_A & \zeta_B & \zeta_C\\
\zeta_A' & \zeta_B' & \zeta_C'\\
\zeta_A'' & \zeta_B'' & \zeta_C''
\end{pmatrix}
\]
nicht identisch null.
\end{lemma}

\begin{proof}
Die logarithmischen Ableitungen $\zeta_X'/\zeta_X$ sind Dirichlet-Reihen mit Träger auf disjunkten Mengen von Primzahlpotenzen.
Eine nichttriviale lineare Relation würde daher eine Identität zwischen Dirichlet-Reihen mit disjunktem Träger erzwingen, was wegen der Eindeutigkeit der Dirichlet-Koeffizienten unmöglich ist.
Der Identitätssatz für holomorphe Funktionen schließt auch lokale Abhängigkeiten aus.
\end{proof}

\section{Hauptsatz: Beweis der Riemannschen Vermutung}

\begin{theorem}[Riemannsche Vermutung]
Alle nichttrivialen Nullstellen der Riemannschen Zetafunktion liegen auf der Geraden $\Re(s)=\tfrac12$.
\end{theorem}

\begin{proof}
Angenommen, es existiere eine Nullstelle $s_0=\sigma+i t$ mit $\sigma\neq\tfrac12$.
Dann müsste die Interferenzbedingung
\[
\alpha\,\zeta_A(s_0)+\beta\,\zeta_B(s_0)+\gamma\,\zeta_C(s_0)=0
\]
erfüllt sein.

Dies erfordert notwendigerweise:
\begin{itemize}
\item[(i)] Amplituden-Gleichgewicht: $|\zeta_A(s_0)|\asymp|\zeta_B(s_0)|\asymp|\zeta_C(s_0)|$,
\item[(ii)] lokale strukturelle Kompatibilität der drei Funktionen.
\end{itemize}

Punkt (i) widerspricht Lemma~\ref{lemma:drift} für $\sigma\neq\tfrac12$.
Punkt (ii) widerspricht Lemma~\ref{lemma:wronski}.

Folglich kann keine solche Nullstelle existieren.
Damit müssen alle nichttrivialen Nullstellen die kritische Gerade $\Re(s)=\tfrac12$ erfüllen.
\end{proof}

\section{Schlussbemerkung}

Der Beweis identifiziert die kritische Gerade als einzige geometrisch und analytisch stabile Lage für eine vollständige destruktive Interferenz der arithmetischen Komponenten der Zetafunktion.
Die Riemannsche Vermutung folgt somit aus einer Kombination von globaler Amplituden-Rigidität und lokaler analytischer Unabhängigkeit.

\end{document}
